\documentclass[11pt]{article}
% shared-preamble.tex
% Common preamble for Paper 1 (paper1-obstruction.tex) and Paper 2 (paper2-recovery.tex).
% Include with: % shared-preamble.tex
% Common preamble for Paper 1 (paper1-obstruction.tex) and Paper 2 (paper2-recovery.tex).
% Include with: % shared-preamble.tex
% Common preamble for Paper 1 (paper1-obstruction.tex) and Paper 2 (paper2-recovery.tex).
% Include with: \input{shared-preamble}

\usepackage[a4paper,margin=1in]{geometry}
\usepackage[T1]{fontenc}
\usepackage[utf8]{inputenc}
\usepackage{lmodern}
\usepackage{amsmath,amssymb,amsthm,mathtools}
\usepackage{mathrsfs}
\usepackage{enumitem}
\usepackage[nameinlink,capitalise,noabbrev]{cleveref}
\usepackage{booktabs}
\usepackage{array}
\usepackage{longtable}
\usepackage{microtype}
\usepackage{xcolor}
\usepackage{float}
\usepackage{graphicx}
\usepackage{url}

% ---------- Theorem environments ----------
\newtheorem{theorem}{Theorem}[section]
\newtheorem{proposition}[theorem]{Proposition}
\newtheorem{lemma}[theorem]{Lemma}
\newtheorem{corollary}[theorem]{Corollary}
\newtheorem{conjecture}[theorem]{Conjecture}
\newtheorem{problem}[theorem]{Problem}
\theoremstyle{definition}
\newtheorem{definition}[theorem]{Definition}
\theoremstyle{remark}
\newtheorem{remark}[theorem]{Remark}

% ---------- cleveref names for custom environments ----------
\crefname{theorem}{Theorem}{Theorems}
\crefname{proposition}{Proposition}{Propositions}
\crefname{lemma}{Lemma}{Lemmas}
\crefname{corollary}{Corollary}{Corollaries}
\crefname{conjecture}{Conjecture}{Conjectures}
\crefname{definition}{Definition}{Definitions}
\crefname{remark}{Remark}{Remarks}

% ---------- Standard math macros ----------
\newcommand{\Q}{\mathbb{Q}}
\newcommand{\Z}{\mathbb{Z}}
\newcommand{\N}{\mathbb{N}}
\newcommand{\C}{\mathbb{C}}
\newcommand{\rank}{\operatorname{rank}}
\newcommand{\Span}{\operatorname{span}}
\newcommand{\Ker}{\operatorname{Ker}}
\newcommand{\SL}{\operatorname{SL}}
\newcommand{\GL}{\operatorname{GL}}

% ---------- Project-specific macros ----------
% Eisenstein-type polynomial span on primes
\newcommand{\Eis}{\mathcal{E}}
% Prime-level cusp obstruction space
\newcommand{\Ob}{\mathcal{O}}
% Cusp sector of a basis
\newcommand{\cusp}{\mathrm{cusp}}
% Prime evaluation map
\newcommand{\ev}{\mathrm{ev}}
% Prime evaluation matrix
\newcommand{\Mmat}{\mathbf{M}}
% Ramanujan tau
\newcommand{\taufn}{\tau}
% Shorthand for MacMahon function
\newcommand{\Ma}[1]{M_{#1}}
% Generating series
\newcommand{\Ua}[1]{U_{#1}}
% Divisor sum
\newcommand{\sig}[1]{\sigma_{#1}}


\usepackage[a4paper,margin=1in]{geometry}
\usepackage[T1]{fontenc}
\usepackage[utf8]{inputenc}
\usepackage{lmodern}
\usepackage{amsmath,amssymb,amsthm,mathtools}
\usepackage{mathrsfs}
\usepackage{enumitem}
\usepackage[nameinlink,capitalise,noabbrev]{cleveref}
\usepackage{booktabs}
\usepackage{array}
\usepackage{longtable}
\usepackage{microtype}
\usepackage{xcolor}
\usepackage{float}
\usepackage{graphicx}
\usepackage{url}

% ---------- Theorem environments ----------
\newtheorem{theorem}{Theorem}[section]
\newtheorem{proposition}[theorem]{Proposition}
\newtheorem{lemma}[theorem]{Lemma}
\newtheorem{corollary}[theorem]{Corollary}
\newtheorem{conjecture}[theorem]{Conjecture}
\newtheorem{problem}[theorem]{Problem}
\theoremstyle{definition}
\newtheorem{definition}[theorem]{Definition}
\theoremstyle{remark}
\newtheorem{remark}[theorem]{Remark}

% ---------- cleveref names for custom environments ----------
\crefname{theorem}{Theorem}{Theorems}
\crefname{proposition}{Proposition}{Propositions}
\crefname{lemma}{Lemma}{Lemmas}
\crefname{corollary}{Corollary}{Corollaries}
\crefname{conjecture}{Conjecture}{Conjectures}
\crefname{definition}{Definition}{Definitions}
\crefname{remark}{Remark}{Remarks}

% ---------- Standard math macros ----------
\newcommand{\Q}{\mathbb{Q}}
\newcommand{\Z}{\mathbb{Z}}
\newcommand{\N}{\mathbb{N}}
\newcommand{\C}{\mathbb{C}}
\newcommand{\rank}{\operatorname{rank}}
\newcommand{\Span}{\operatorname{span}}
\newcommand{\Ker}{\operatorname{Ker}}
\newcommand{\SL}{\operatorname{SL}}
\newcommand{\GL}{\operatorname{GL}}

% ---------- Project-specific macros ----------
% Eisenstein-type polynomial span on primes
\newcommand{\Eis}{\mathcal{E}}
% Prime-level cusp obstruction space
\newcommand{\Ob}{\mathcal{O}}
% Cusp sector of a basis
\newcommand{\cusp}{\mathrm{cusp}}
% Prime evaluation map
\newcommand{\ev}{\mathrm{ev}}
% Prime evaluation matrix
\newcommand{\Mmat}{\mathbf{M}}
% Ramanujan tau
\newcommand{\taufn}{\tau}
% Shorthand for MacMahon function
\newcommand{\Ma}[1]{M_{#1}}
% Generating series
\newcommand{\Ua}[1]{U_{#1}}
% Divisor sum
\newcommand{\sig}[1]{\sigma_{#1}}


\usepackage[a4paper,margin=1in]{geometry}
\usepackage[T1]{fontenc}
\usepackage[utf8]{inputenc}
\usepackage{lmodern}
\usepackage{amsmath,amssymb,amsthm,mathtools}
\usepackage{mathrsfs}
\usepackage{enumitem}
\usepackage[nameinlink,capitalise,noabbrev]{cleveref}
\usepackage{booktabs}
\usepackage{array}
\usepackage{longtable}
\usepackage{microtype}
\usepackage{xcolor}
\usepackage{float}
\usepackage{graphicx}
\usepackage{url}

% ---------- Theorem environments ----------
\newtheorem{theorem}{Theorem}[section]
\newtheorem{proposition}[theorem]{Proposition}
\newtheorem{lemma}[theorem]{Lemma}
\newtheorem{corollary}[theorem]{Corollary}
\newtheorem{conjecture}[theorem]{Conjecture}
\newtheorem{problem}[theorem]{Problem}
\theoremstyle{definition}
\newtheorem{definition}[theorem]{Definition}
\theoremstyle{remark}
\newtheorem{remark}[theorem]{Remark}

% ---------- cleveref names for custom environments ----------
\crefname{theorem}{Theorem}{Theorems}
\crefname{proposition}{Proposition}{Propositions}
\crefname{lemma}{Lemma}{Lemmas}
\crefname{corollary}{Corollary}{Corollaries}
\crefname{conjecture}{Conjecture}{Conjectures}
\crefname{definition}{Definition}{Definitions}
\crefname{remark}{Remark}{Remarks}

% ---------- Standard math macros ----------
\newcommand{\Q}{\mathbb{Q}}
\newcommand{\Z}{\mathbb{Z}}
\newcommand{\N}{\mathbb{N}}
\newcommand{\C}{\mathbb{C}}
\newcommand{\rank}{\operatorname{rank}}
\newcommand{\Span}{\operatorname{span}}
\newcommand{\Ker}{\operatorname{Ker}}
\newcommand{\SL}{\operatorname{SL}}
\newcommand{\GL}{\operatorname{GL}}

% ---------- Project-specific macros ----------
% Eisenstein-type polynomial span on primes
\newcommand{\Eis}{\mathcal{E}}
% Prime-level cusp obstruction space
\newcommand{\Ob}{\mathcal{O}}
% Cusp sector of a basis
\newcommand{\cusp}{\mathrm{cusp}}
% Prime evaluation map
\newcommand{\ev}{\mathrm{ev}}
% Prime evaluation matrix
\newcommand{\Mmat}{\mathbf{M}}
% Ramanujan tau
\newcommand{\taufn}{\tau}
% Shorthand for MacMahon function
\newcommand{\Ma}[1]{M_{#1}}
% Generating series
\newcommand{\Ua}[1]{U_{#1}}
% Divisor sum
\newcommand{\sig}[1]{\sigma_{#1}}


\title{The Weight-12 Cusp Obstruction in\\
       Partition-Theoretic Prime Detection}
\author{Nigel Randsley}
\date{}

\begin{document}

\maketitle

\begin{abstract}
We prove that the weight-12 cusp form $\Delta$ forces a structural obstruction in the
program of detecting primes by polynomial combinations of MacMahon partition functions.
Specifically, for any polynomial degree $d\ge 0$, no non-trivial prime-vanishing expression
formed from the basis $\{\Ma{1},\dots,\Ma{6}\}$ can involve $\Ma{6}$.
The proof rests on three lemmas: (i) the prime restrictions of $\Ma{1},\dots,\Ma{5}$ lie
entirely in the polynomial ring $\Q[p]$; (ii) the prime restriction of $\Ma{6}$ contains a
non-zero cusp contribution $c\cdot\taufn(p)$ with $c = -17/150\,450\,048\,000$;
(iii) no nonzero polynomial multiple of $\taufn(p)$ is a polynomial function of $p$
on the primes, a consequence of the Sato-Tate equidistribution theorem and the Deligne bound.
The theorem converts a computational observation of the companion paper into a structural
result and identifies the Bernoulli prime $691$ as the arithmetic signature of the
obstruction.
We pose the question of whether analogous obstructions appear at each weight where the
cusp-form dimension first increases.
\end{abstract}

\tableofcontents

% ==========================================================================
\section{Introduction}
\label{sec:intro}
% ==========================================================================

\subsection{Context}

Craig, van Ittersum, and Ono \cite{CIO2024} showed that suitable non-negative polynomial
combinations of MacMahon partition functions $\Ma{a}(n)$ vanish at the primes and only at
the primes, producing four explicit expressions $E_1,E_2,E_3,E_4$.
The generating series $\Ua{a}(q)=\sum_{n\ge 1}\Ma{a}(n)q^n$ is a quasimodular form of
weight $2a$ on $\SL_2(\Z)$ in the sense of Kaneko--Zagier \cite{KZ1995}.
For $a\le 5$ (weight $\le 10$), these spaces are spanned entirely by Eisenstein series,
so $\Ma{a}(p)$ is a polynomial in $p$ at every prime $p$.
The expressions $E_1,\dots,E_4$ arise from the resulting polynomial dependencies.

Van Ittersum, Mauth, Ono, and Singh \cite{VIMOS2025} and Kane, Krishnamoorthy, and Lau
\cite{KKL2025} independently proved the CIO conjecture: every level-1 prime-detecting
quasimodular form is an Eisenstein series.
This closes the story at weights $\le 10$, but leaves open what happens when cusp forms
enter.

The companion paper \cite{Randsley2025b} establishes computationally that:
\begin{itemize}[leftmargin=*]
  \item Within the basis $\{\Ma{1},\dots,\Ma{6}\}$, no new prime-vanishing direction involving
        $\Ma{6}$ appears at any tested polynomial degree $d\in\{0,\dots,6\}$.
  \item A new direction $E_5$ does appear at degree 2 in the restricted basis
        $\{\Ma{1},\dots,\Ma{5}\}$.
  \item A further direction $E_6$ appears at degree 4 only after $\Ma{7}$ is adjoined.
\end{itemize}
The present paper proves the first of these observations as a theorem.

\subsection{Main Results}

\begin{theorem}[Obstruction Theorem]
\label{thm:obstruction}
For every polynomial degree $d\ge 0$, the prime evaluation map
\[
  \ev_{\mathrm{primes}} \colon
  V_d^{(6)} \longrightarrow \Q^{\infty},
  \quad
  V_d^{(6)} = \Span_\Q\bigl\{n^k \Ma{a}(n) : 0\le k\le d,\; 1\le a\le 6\bigr\},
\]
has kernel equal to the kernel of its restriction to $V_d^{(5)}$.
In particular,
\[
  \dim\Ker\bigl(\ev_{\mathrm{primes}}\big|_{V_d^{(6)}}\bigr)
  = \dim\Ker\bigl(\ev_{\mathrm{primes}}\big|_{V_d^{(5)}}\bigr).
\]
\end{theorem}

\begin{corollary}
\label{cor:M6-absent}
No prime-vanishing expression formed from $\{\Ma{1},\dots,\Ma{6}\}$ at any polynomial
degree can have a non-zero $\Ma{6}$-coefficient.
\end{corollary}

\begin{remark}[Why $E_5$ still exists]
Corollary~\ref{cor:M6-absent} rules out $\Ma{6}$ within the basis $\{\Ma{1},\dots,\Ma{6}\}$,
but does not preclude new prime-vanishing directions that avoid $\Ma{6}$ entirely.
Indeed, the companion computational paper \cite{Randsley2025b} finds a new direction $E_5$
at degree 2 within the restricted basis $\{\Ma{1},\dots,\Ma{5}\}$.
Since these functions have only Eisenstein prime restrictions, the linear dependencies
among $\{p^k\Ma{a}(p) : 0\le k\le 2, 1\le a\le 5\}$ are purely polynomial, and a new
null direction can appear once enough polynomial degree is available.
The weight-12 obstruction is specifically an obstruction to any expression \emph{involving}
$\Ma{6}$; it leaves the $\{\Ma{1},\dots,\Ma{5}\}$ subspace unaffected.
\end{remark}

The proof of \cref{thm:obstruction} occupies \cref{sec:M6-prime,sec:independence}.
\Cref{sec:arithmetic} interprets the arithmetic signature of the obstruction.
\Cref{sec:computational} records the supporting computational data.

\paragraph{Proof roadmap.}
The proof has a short and rigid structure.
First, the lower MacMahon functions $M_1,\dots,M_5$ are shown to have polynomial prime
restrictions, so the restricted prime-evaluation space is Eisenstein-type.
Second, $M_6$ is written in the form
\[
  M_6(p) = f(p) + c_\tau\tau(p),
\]
with $f\in\mathbb{Q}[p]$ and $c_\tau\ne 0$.
Third, a non-cancellation lemma shows that no nonzero polynomial multiple of $\tau(p)$
can lie in the Eisenstein-type prime span.
The obstruction theorem then follows by decomposing any candidate prime-vanishing relation
into its polynomial part and its cusp part and observing that the latter cannot vanish
unless the entire $M_6$-sector is zero.

% ==========================================================================
\section{Background}
\label{sec:background}
% ==========================================================================

\subsection{MacMahon Partition Functions}

For $a\ge 1$ define
\[
  \Ma{a}(n) = \sum_{\substack{0<s_1<\cdots<s_a \\ m_i\ge 1,\;\sum m_i s_i=n}}
              m_1 m_2\cdots m_a.
\]
The generating series $\Ua{a}(q)=\sum_{n\ge 1}\Ma{a}(n)q^n$ is a quasimodular form of
weight $2a$ and depth at most $a-1$ on $\SL_2(\Z)$ \cite{CIO2024,KZ1995}.

\subsection{Prime Evaluation and Divisor Sums}

For $a\le 5$, the quasimodular space of weight $2a$ is spanned by Eisenstein series, so
$\Ma{a}(n)$ is a polynomial combination of divisor power sums
$\sig{k}(n)=\sum_{d\mid n}d^k$.
At a prime $p$, the divisor sum evaluates as
\[
  \sig{k}(p) = p^k + 1 \in \Q[p],
\]
so $\Ma{a}(p)\in\Q[p]$ for every $a\le 5$.

\subsection{Modular Forms at Weight 12}

The space of modular forms of weight 12 on $\SL_2(\Z)$ is two-dimensional:
\[
  M_{12}(\SL_2(\Z)) = \Q\cdot E_{12} \oplus \Q\cdot\Delta,
\]
where $E_{12}$ is the Eisenstein series of weight 12 and
$\Delta(q) = q\prod_{k=1}^{\infty}(1-q^k)^{24} = \sum_{n=1}^{\infty}\taufn(n)q^n$
is the unique normalized cusp form, with $\taufn$ denoting Ramanujan's $\tau$-function.
Since $\dim S_{12}(\SL_2(\Z))=1$ and $\Ua{6}$ has non-trivial projection onto this
cusp-form space, the closed form of $\Ma{6}$ necessarily involves $\taufn$.

The $\tau$-function satisfies the Deligne bound $|\taufn(p)|\le 2p^{11/2}$ for every prime
$p$ \cite{Deligne1974}.
Crucially, $\taufn(p)$ is \emph{not} a polynomial in $p$; establishing this rigorously is
the crux of \cref{sec:independence}.

% ==========================================================================
\section{The Closed Form of $M_6$ and Its Prime Restriction}
\label{sec:M6-prime}
% ==========================================================================

\begin{lemma}[Eisenstein-type span for $a\le 5$]
\label{lem:M1-M5-poly}
For $1\le a\le 5$, the function $p\mapsto\Ma{a}(p)$ is a polynomial in $p$ with rational
coefficients.
More precisely, $\Ma{a}(p)\in\Q[p]$ for all primes $p$.
\end{lemma}

\begin{proof}
For $a\le 5$, the quasimodular weight $2a\le 10$, and the relevant spaces are generated by
Eisenstein series.
Hence $\Ma{a}(n)$ can be written as a $\Q$-linear combination of monomials
$n^j\sig{k}(n)$.
Evaluating at a prime $p$: since $p$ is prime, the only divisors of $p$ are $1$ and $p$,
so $\sig{k}(p)=1+p^k$.
Substituting, every $n^j\sig{k}(n)$ evaluates to a polynomial in $p$.
\end{proof}

Let $\Eis$ denote the $\Q$-algebra of functions $f\colon\{\text{primes}\}\to\Q$ of the
form $p\mapsto g(p)$ for some $g\in\Q[p]$.
Lemma~\ref{lem:M1-M5-poly} above states that $\Ma{a}\big|_{\mathrm{primes}}\in\Eis$ for $a\le 5$.

\begin{lemma}[Cusp contribution of $M_6$]
\label{lem:M6-cusp}
The prime restriction of $\Ma{6}$ takes the form
\[
  \Ma{6}(p) = f(p) + c_\tau\cdot\taufn(p),
\]
where $f\in\Q[p]$ and
\[
  c_\tau = -\frac{17}{150\,450\,048\,000} \ne 0.
\]
\end{lemma}

\begin{proof}
Fitting a $55\times 22$ rational linear system (see \cref{app:M6-formula}) yields the
exact closed form
\[
  \Ma{6}(n) = \sum_{j=0}^{5} P_j(n)\,\sig{2j+1}(n)
              - \frac{17}{150\,450\,048\,000}\,\taufn(n),
\]
with $P_j\in\Q[n]$ explicitly computed (see \cref{app:computations}).
The system has 55 equations (one per data point $n=1,\dots,55$) and 22 unknowns
(six polynomial coefficients per $P_j$ plus the $\taufn$-coefficient $c_\tau$);
since the $55\times 22$ matrix has full column rank 22 over $\Q$ (verified by exact
rational RREF), the solution is unique.
The fitted formula is used here only as an exact coefficient identity: once the ansatz
space is fixed and the corresponding rational system is shown to have full column rank,
the recovered coefficients are uniquely determined and may be used in the subsequent proof
without heuristic ambiguity.
Evaluating at a prime $p$: each $\sig{2j+1}(p)=p^{2j+1}+1$, so the divisor-sum part
$\sum_j P_j(p)\,\sig{2j+1}(p)\in\Q[p]$.
The $\taufn$-part contributes $c_\tau\cdot\taufn(p)$, which is not in $\Eis$ by
Lemma~\ref{thm:tau-independence} below.
Since $c_\tau=-17/150\,450\,048\,000\ne 0$, the cusp contribution is non-trivial.
\end{proof}

% ==========================================================================
\section{Linear Independence of $\tau$ over Polynomials at Primes}
\label{sec:independence}
% ==========================================================================

\begin{lemma}[Non-polynomiality of $\tau$ at primes]
\label{thm:tau-independence}
The function $p\mapsto\taufn(p)$ does not lie in $\Eis$: there is no polynomial
$g\in\Q[x]$ such that $\taufn(p)=g(p)$ for all sufficiently large primes $p$.
\end{lemma}

\begin{proof}
This proof uses the Sato-Tate theorem \cite{BLGHT2011}, which we state as a black box:
the normalized values $\taufn(p)/(2p^{11/2})$ are equidistributed on $[-1,1]$ with respect
to the measure $\mu_{\mathrm{ST}}=\frac{2}{\pi}\sqrt{1-t^2}\,dt$ as $p\to\infty$.

Suppose for contradiction that $\taufn(p)=g(p)$ for all sufficiently large primes $p$,
where $g\in\Q[x]$ is a polynomial of degree $m\ge 0$.
The Deligne bound \cite{Deligne1974} gives $|\taufn(p)|\le 2p^{11/2}$, so in particular
$|g(p)|\le 2p^{11/2}$ for large $p$, forcing $m\le 11/2$, i.e.\ $m\le 5$.

Now consider the sign of $\taufn(p)$.
By Sato-Tate equidistribution, the set of primes $p$ for which $\taufn(p)>0$ has density
$\mu_{\mathrm{ST}}((0,1]) = 1/2$, and the set with $\taufn(p)<0$ also has density $1/2$.
In particular, $\taufn(p)$ changes sign on a positive-density set of primes.
However, a nonzero polynomial $g$ of degree $m$ has at most $m$ real roots, hence is
eventually of constant sign.
Since $m\le 5$ is finite, $g(p)$ is eventually positive or eventually negative, so $g$
cannot equal $\taufn(p)$ on any positive-density set of primes.
This is the required contradiction.
\end{proof}

\begin{remark}[Dependence on Sato-Tate]
Lemma~\ref{thm:tau-independence} is proved using the Sato-Tate theorem, established by
Barnet-Lamb, Geraghty, Harris, and Taylor \cite{BLGHT2011}.
The conclusion also follows from Serre's open image theorem \cite{Serre1972} combined
with the Chebotarev density theorem via a standard Galois-theoretic argument, which
yields the sign-change conclusion without appealing to equidistribution.
Both proofs use deep results from the theory of $\ell$-adic representations; no fully
elementary proof is currently known.
\end{remark}

\begin{lemma}[Polynomial multiples of $\tau$ cannot be polynomial]
\label{lem:no-poly-combo}
Let $Q\in\Q[x]$ be a nonzero polynomial.
Then the function $p\mapsto Q(p)\cdot\taufn(p)$ is not identically zero on the primes,
and in particular $Q(p)\cdot\taufn(p)\notin\Eis$.
\end{lemma}

\begin{proof}
Suppose $Q(p)\cdot\taufn(p) = 0$ for all sufficiently large primes $p$.
Since $Q$ is a nonzero polynomial, it has only finitely many real roots, so $Q(p)\ne 0$
for all but finitely many primes.
Therefore $\taufn(p)=0$ for all but finitely many primes.

However, by the Sato-Tate theorem \cite{BLGHT2011}, the normalized values
$\taufn(p)/(2p^{11/2})$ are equidistributed on $[-1,1]$ under $\mu_{\mathrm{ST}}$.
In particular, the density of primes where $\taufn(p)>0$ is $\mu_{\mathrm{ST}}((0,1])=1/2>0$.
Having $\taufn(p)=0$ for all but finitely many primes would force the density of
$\{\taufn(p)>0\}$ to be zero, contradicting this.

Therefore $Q(p)\cdot\taufn(p)\ne 0$ for infinitely many primes, i.e.\ the product is not
identically zero.
For the second claim: suppose $Q(p)\cdot\taufn(p)=g(p)$ for some $g\in\Q[p]$ and all
sufficiently large primes $p$.
Since $Q\in\Q[x]$ is nonzero, it has only finitely many real roots, so $Q(p)\ne 0$ for
all but finitely many primes; at such primes, $\taufn(p) = g(p)/Q(p)$.
Let $h = g/Q\in\Q(x)$ be the rational function.
Since $Q(p)\ne 0$ for all sufficiently large primes, $h$ has no poles at sufficiently
large primes, and $\taufn(p) = h(p)$ for all sufficiently large primes.

Now, a nonzero rational function $h\in\Q(x)$ has eventually constant sign: writing
$h(x) = c\,x^m(1+O(1/x))$ for large $x$ (where $c\ne 0$ and $m$ is the
degree difference of numerator and denominator), the sign of $h(x)$ equals the sign of
$c$ for all sufficiently large $x$.
Therefore $h(p)$ is eventually positive or eventually negative on the primes.
However, the Sato-Tate theorem \cite{BLGHT2011} gives $\taufn(p)>0$ for a positive-density
set of primes (density $1/2$) and $\taufn(p)<0$ for another set of density $1/2$, so
$\taufn(p)$ does not have eventually constant sign — contradiction.
\end{proof}

% ==========================================================================
\section{Proof of the Obstruction Theorem}
\label{sec:proof}
% ==========================================================================

\begin{proof}[Proof of \cref{thm:obstruction}]
Let $F\in V_d^{(6)}$ be a prime-vanishing expression:
\[
  F(n) = \sum_{k=0}^{d}\sum_{a=1}^{6} c_{k,a}\,n^k\Ma{a}(n),
  \quad c_{k,a}\in\Q,
\]
with $F(p)=0$ for every prime $p$.
We show $c_{k,6}=0$ for all $k$.

Evaluating at a prime $p$ and separating the $a=6$ terms:
\[
  \underbrace{\sum_{k=0}^{d}\sum_{a=1}^{5} c_{k,a}\,p^k\Ma{a}(p)}_{=:\,A(p)}
  +
  \underbrace{\sum_{k=0}^{d} c_{k,6}\,p^k}_{=:\,Q(p)}
  \cdot\Ma{6}(p)
  = 0.
\]

By Lemma~\ref{lem:M1-M5-poly}, each $\Ma{a}(p)\in\Q[p]$ for $a\le 5$, so $A(p)\in\Q[p]$.

By Lemma~\ref{lem:M6-cusp}, $\Ma{6}(p)=f(p)+c_\tau\taufn(p)$ with $f\in\Q[p]$ and
$c_\tau\ne 0$.
Substituting:
\[
  A(p) + Q(p)f(p) + c_\tau Q(p)\taufn(p) = 0.
\]

The sum $A(p)+Q(p)f(p)\in\Q[p]$ (a polynomial in $p$).
Rearranging: $c_\tau Q(p)\cdot\taufn(p) = -(A(p)+Q(p)f(p))\in\Q[p]$.
If $Q\not\equiv 0$, then $c_\tau Q\in\Q[p]$ is a nonzero polynomial, and
Lemma~\ref{lem:no-poly-combo} gives $c_\tau Q(p)\cdot\taufn(p)\notin\Eis$, contradicting
the right-hand side being a polynomial.
Therefore $Q(p)\equiv 0$, i.e.\ $c_{k,6}=0$ for all $k$.

With all $M_6$-coefficients zero, $F\in V_d^{(5)}$ and the claim follows.
\end{proof}

\begin{proof}[Proof of Corollary~\ref{cor:M6-absent}]
Immediate from \cref{thm:obstruction}: any prime-vanishing element of $V_d^{(6)}$ lies in
$V_d^{(5)}$, so its $\Ma{6}$-coefficient is zero.
\end{proof}

% ==========================================================================
\section{Arithmetic of the Obstruction}
\label{sec:arithmetic}
% ==========================================================================

The coefficient $c_\tau = -17/150\,450\,048\,000$ has denominator
\[
  150\,450\,048\,000 = 2^{10}\times 3^5\times 5^3\times 7\times 691.
\]
The prime $691$ arises as the numerator of the Bernoulli number $B_{12}$:
$B_{12}=-691/2730$ \cite{Bernoulli}.
The Ramanujan congruence $\taufn(n)\equiv\sig{11}(n)\pmod{691}$ reflects the fact that
$E_{12}$ and $\Delta$ are congruent modulo 691.
In particular, the weight-12 obstruction is \emph{arithmetically marked} by 691.

\begin{remark}
The same factor 691 appears in the leading $\Ma{7}$-coefficient of $E_6$:
$892\,705\,701\,888\,000 = 2^{10}\times 3^5\times 5^3\times 7\times 691\times 1303$,
suggesting that the recovery mechanism (Paper 2, \cite{Randsley2025c}) also carries
this arithmetic signature.
\end{remark}

% ==========================================================================
\section{Computational Verification}
\label{sec:computational}
% ==========================================================================

The proof of \cref{thm:obstruction} is unconditional, but the companion computation
provides concrete evidence and motivates the structural result.

\begin{proposition}[Computational confirmation]
For each tested degree $d\in\{0,1,2,3,4,5,6\}$ and $N=95$ primes in $[2,500]$,
\[
  \rank\,\Mmat(d,6,95) = \rank\,\Mmat(d,5,95) + (d+1),
\]
confirming that all $\Ma{6}$-columns are pivot columns and the null space gains no new
direction.
\end{proposition}

See \cref{app:rank-data} for the rank and nullity data.

% ==========================================================================
\section{Open Questions}
\label{sec:open}
% ==========================================================================

\begin{enumerate}[leftmargin=*,label=\textbf{(\arabic*)}]

\item \textbf{Higher-weight obstructions.}
Does an analogous obstruction arise at every weight $2a$ for which $\dim S_{2a}$ strictly
increases?
The first such increases occur at weights 12, 16, 18, 20, 22, 24 (where $\dim S_{24}=2$).
The companion paper \cite{Randsley2025b} finds no new $E_i$ through $a_{\max}=12$, but
a structural theorem predicting obstruction at each new cusp-dimension layer remains open.

\item \textbf{Higher levels.}
The CIO framework has been extended to level $N$ by Kang, Matsusaka, and Shin
\cite{KMS2024}.
Does the weight-12 obstruction persist, weaken, or acquire new structure at higher levels?

\item \textbf{Effective non-polynomiality.}
Lemma~\ref{thm:tau-independence} is proved via the Sato-Tate theorem and open image.
Can one give an explicit, elementary proof using only finite prime computations and
Ramanujan's congruences?

\end{enumerate}

% ==========================================================================
\section{Conclusion}
\label{sec:conclusion}
% ==========================================================================

We have proved that the weight-12 cusp form $\Delta$ creates a structural obstruction:
the $\tau$-contribution to $\Ma{6}(p)$ is algebraically independent of all polynomial
functions of $p$, so no prime-vanishing expression in the basis
$\{\Ma{1},\dots,\Ma{6}\}$ can involve $\Ma{6}$.
The obstruction is arithmetically signed by the Bernoulli prime $691$ and is specific
to the restricted basis; the companion paper \cite{Randsley2025c} shows it can be
resolved by adjoining $\Ma{7}$.

\medskip
\noindent\textbf{Acknowledgements.}
The computations supporting this paper were performed in exact rational arithmetic
using Julia scripts available at \url{https://github.com/randsley/PartitionPrimes}.
The author thanks colleagues for helpful discussions (names to be supplied).

% ==========================================================================
\begin{thebibliography}{99}

% ---- Primary source ----
\bibitem{CIO2024}
W.~Craig, J.-W.~van Ittersum, and K.~Ono,
\emph{Integer partitions detect the primes},
Proc.\ Natl.\ Acad.\ Sci.\ USA \textbf{121} (2024), no.~39, e2409417121.
\newblock arXiv:2405.06451.

% ---- The Eisenstein conjecture ----
\bibitem{VIMOS2025}
J.-W.~van Ittersum, L.~Mauth, K.~Ono, and A.~Singh,
\emph{Quasimodular forms that detect primes are Eisenstein},
preprint (2025).
\newblock arXiv:2507.20432.

\bibitem{KKL2025}
B.~Kane, K.~Krishnamoorthy, and Y.-K.~Lau,
\emph{Prime detecting quasi-modular forms in higher level},
preprint (2025).
\newblock arXiv:2511.04030.

% ---- Level-N extensions ----
\bibitem{KMS2024}
S.-Y.~Kang, T.~Matsusaka, and G.~Shin,
\emph{Quasi-modularity in MacMahon partition variants and prime detection},
Ramanujan J.\ \textbf{67} (2025), no.~4, Art.~81.
\newblock arXiv:2412.19180.

% ---- Companion papers ----
\bibitem{Randsley2025b}
N.~Randsley,
\emph{Extending Partition-Theoretic Prime Detection: A Computational Study of the
Weight-12 Barrier, the Expression $E_5$, and a Basis-Dependent Recovery of $E_6$},
preprint (2025).

\bibitem{Randsley2025c}
N.~Randsley,
\emph{Recovery of a Prime-Vanishing Direction via $M_6$ and $M_7$:
Structural Proof and Obstruction Dimension},
preprint (2025).

% ---- Background: quasimodular forms ----
\bibitem{KZ1995}
M.~Kaneko and D.~Zagier,
\emph{A generalized Jacobi theta function and quasimodular forms},
in \textit{The Moduli Space of Curves}
  (R.~Dijkgraaf, C.~Faber, and G.~van~der~Geer, eds.),
Progr.\ Math.\ \textbf{129}, Birkh\"auser, Boston, 1995, pp.~165--172.

% ---- Background: Weil conjectures and Deligne bounds ----
\bibitem{Deligne1974}
P.~Deligne,
\emph{La conjecture de Weil,~I},
Inst.\ Hautes \'Etudes Sci.\ Publ.\ Math.\ \textbf{43} (1974), 273--307.

% ---- Background: Sato-Tate theorem ----
\bibitem{BLGHT2011}
T.~Barnet-Lamb, D.~Geraghty, M.~Harris, and R.~Taylor,
\emph{A family of Calabi--Yau varieties and potential automorphy~II},
Publ.\ Res.\ Inst.\ Math.\ Sci.\ \textbf{47} (2011), 29--98.

% ---- Background: open image theorem ----
\bibitem{Serre1972}
J.-P.~Serre,
\emph{Propri\'et\'es galoisiennes des points d'ordre fini des courbes elliptiques},
Invent.\ Math.\ \textbf{15} (1972), 259--331.

% ---- Background: Bernoulli numbers and modular forms ----
\bibitem{Bernoulli}
H.~Cohen and F.~Str\"omberg,
\emph{Modular Forms: A Classical and Computational Introduction},
Graduate Studies in Mathematics, vol.~179, American Mathematical Society,
Providence, 2017.

\end{thebibliography}

% ==========================================================================
% appendix-computations.tex
% Shared computational appendix for Paper 1 and Paper 2.
% Include with: % appendix-computations.tex
% Shared computational appendix for Paper 1 and Paper 2.
% Include with: % appendix-computations.tex
% Shared computational appendix for Paper 1 and Paper 2.
% Include with: \input{appendix-computations}

\section*{Appendix: Computational Data}
\addcontentsline{toc}{section}{Appendix: Computational Data}
\label{app:computations}

All computations were performed in exact rational arithmetic.
Prime evaluation matrices were built from the first $N$ primes and reduced via
rational row-reduction (RREF over $\Q$).
The source code is available at the project repository.

% ------------------------------------------------------------------
\subsection*{A.1\quad Closed Form of $M_6$}
\label{app:M6-formula}

Fitting the $55\times 22$ rational linear system
(21 polynomial-weighted divisor-sum columns plus one $\tau$-column)
yields the exact formula
\[
  M_6(n) = \sum_{j=0}^{5} P_j(n)\,\sigma_{2j+1}(n)
            - \frac{17}{150\,450\,048\,000}\,\tau(n),
\]
with explicit polynomials $P_j\in\Q[n]$:
\begin{align*}
P_0(n)&=\frac{550499}{4541644800}-\frac{153617}{371589120}\,n+\frac{2159}{6635520}\,n^2-\frac{67}{737280}\,n^3+\frac{11}{1105920}\,n^4-\frac{1}{2764800}\,n^5\\[4pt]
P_1(n)&=\frac{153617}{743178240}-\frac{2159}{8847360}\,n+\frac{67}{819200}\,n^2-\frac{11}{1105920}\,n^3+\frac{1}{2580480}\,n^4\\[4pt]
P_2(n)&=\frac{2159}{88473600}-\frac{67}{4915200}\,n+\frac{11}{5160960}\,n^2-\frac{1}{10321920}\,n^3\\[4pt]
P_3(n)&=\frac{67}{123863040}-\frac{11}{74317824}\,n+\frac{1}{111476736}\,n^2\\[4pt]
P_4(n)&=\frac{11}{3715891200}-\frac{1}{3096576000}\,n\\[4pt]
P_5(n)&=\frac{1}{268995133440}
\end{align*}
These 22 coefficients were verified against direct partition enumeration for
$n=1,\ldots,30$ with no mismatches (see repository script
\texttt{Paper/compute\_m6\_coefficients.jl}).
The key datum is the $\tau$-coefficient
\[
  c_\tau = -\frac{17}{150\,450\,048\,000},
  \qquad
  150\,450\,048\,000 = 2^{10}\times 3^5\times 5^3\times 7\times 691.
\]
The denominator prime $691 = \mathrm{num}(B_{12})$ is the Bernoulli prime
indexing the weight-12 Eisenstein series.

% ------------------------------------------------------------------
\subsection*{A.2\quad Cusp Structure of $M_7$ and Value of $\beta$}
\label{app:M7-formula}

Despite $\dim S_{14}(\SL_2(\Z))=0$, the closed form of $M_7(n)$ involves
$n\cdot\tau(n)$.
This follows from the quasimodular decomposition of $U_7(q)$: the depth-1 component
lies in $E_2\cdot M_{12}$, and the identity $D\Delta = E_2\Delta$ (proved from Ramanujan's differential equations
in the companion paper) shows that
$[q^n](E_2\Delta)=n\cdot\tau(n)$ exactly.
There is no independent $\tau(n)$ term: the $\tau(p)$ direction in $\mathcal{O}_d$
comes entirely from $M_6$ (via $c_\tau\ne 0$).

The exact closed form is
\[
  M_7(n) = \sum_{j=0}^{6} Q_j(n)\,\sigma_{2j+1}(n) + \beta\,n\cdot\tau(n),
\]
with $\beta\ne 0$.
The exact rational value of $\beta$ is the coefficient of $E_2\Delta$ in the
quasimodular decomposition of $U_7$, computed in exact arithmetic by the repository
script \texttt{Paper/extract\_e6\_e7.jl} (function \texttt{\_fit\_m7!}, column~30
of the augmented RREF system) and verified against direct partition enumeration for
$n=1,\dots,30$.

The structural significance is:
\begin{itemize}
  \item $\beta\ne 0$: $M_7$ contributes $p\cdot\tau(p)$ at primes, a direction absent
        from $M_6$ alone (which contributes only $\tau(p)$).
\end{itemize}
Together, the $\{M_6,M_7\}$ sub-system at primes carries both independent cusp
directions $\tau(p)$ and $p\cdot\tau(p)$, which is the source of the
cusp-cancellation mechanism underlying $E_6$.

% ------------------------------------------------------------------
\subsection*{A.3\quad Prime Evaluation Rank Data}
\label{app:rank-data}

\begin{table}[H]
\centering
\caption{Rank and nullity of prime evaluation matrices at $N=95$ primes,
         showing the pivot status of $M_6$ within the restricted basis.}
\label{tab:app-pivot}
\begin{tabular}{@{}lcccc@{}}
\toprule
Matrix & Rows & Cols & Rank & Nullity \\
\midrule
$a_{\max}=5,\;d=2$ & 95 & 15 & 11 & 4 \\
$a_{\max}=6,\;d=2$ & 95 & 18 & 14 & 4 \\
$a_{\max}=5,\;d=3$ & 95 & 20 & 12 & 8 \\
$a_{\max}=6,\;d=3$ & 95 & 24 & 16 & 8 \\
\bottomrule
\end{tabular}
\end{table}

In each case the nullity does not increase when $M_6$ is adjoined,
confirming that all $M_6$-columns are pivot columns.

% ------------------------------------------------------------------
\subsection*{A.4\quad Nullity Gains When $M_7$ Is Adjoined}
\label{app:nullity-m7}

\begin{table}[H]
\centering
\caption{Nullity comparison at $N=95$ primes.
         The gain first appears at $d=4$.}
\label{tab:app-nullity-m7}
\begin{tabular}{@{}ccccc@{}}
\toprule
$d$ & Nullity $\{M_1\text{--}M_5\}$ &
      Nullity $\{M_1\text{--}M_6\}$ &
      Nullity $\{M_1\text{--}M_7\}$ & Gain \\
\midrule
2 & 4  & 4  & 4  & 0 \\
3 & 8  & 8  & 8  & 0 \\
4 & 12 & 12 & 13 & $+1$ \\
5 & 16 & 16 & 18 & $+2$ \\
6 & 20 & 20 & 23 & $+3$ \\
\bottomrule
\end{tabular}
\end{table}

% ------------------------------------------------------------------
\subsection*{A.5\quad Explicit Formula for $E_5$}
\label{app:E5-formula}

After normalization to primitive integer coefficients,
the unique new prime-vanishing direction at degree~2 in the basis
$\{M_1,\dots,M_5\}$ is
\begin{align*}
E_5(n) ={}
  & (-450450+675675n-225225n^2)\,M_1(n) \\
  &+(960960n-120120n^2)\,M_2(n) \\
  &+(2534912n-166016n^2)\,M_3(n) \\
  &+(7999488n-322560n^2)\,M_4(n) \\
  &+258048000\,M_5(n).
\end{align*}
The constant $M_5$-coefficient is
$258048000 = 2^{11}\times 3^2\times 5^3\times 7\times 127$.

% ------------------------------------------------------------------
\subsection*{A.6\quad Explicit Formula for $E_6$}
\label{app:E6-formula}

The unique new prime-vanishing direction at degree~4 in the basis
$\{M_1,\dots,M_7\}$, normalized analogously, is
\begin{align*}
E_6(n) ={}
  &(105367732470 - 277943208789n + 289613157079n^2 \\
  &\quad{}-145583618619n^3+28545937859n^4)\,M_1(n)\\
  &+(-51324522800n^3+4292382480n^4)\,M_2(n)\\
  &+(-40313554176n^3+1776519808n^4)\,M_3(n)\\
  &+(-32888346624n^3+770273280n^4)\,M_4(n)\\
  &+(-7741440000n^3-154828800n^4)\,M_5(n)\\
  &+(-483548921856000+37196070912000n)\,M_6(n)\\
  &+892705701888000\,M_7(n).
\end{align*}
The leading $M_7$-coefficient is
$892705701888000 = 2^{10}\times 3^5\times 5^3\times 7\times 691\times 1303$.
The factor $691$ connects $E_6$ arithmetically to the weight-12 modular structure of
$M_6$.

% ------------------------------------------------------------------
\subsection*{A.7\quad Extended Search to $a_{\max}=12$}
\label{app:extended-search}

\begin{table}[H]
\centering
\caption{Extended search to $a_{\max}=12$.
         No independent $E_i$ was found beyond $E_6$.}
\label{tab:app-extended}
\begin{tabular}{@{}ccccc@{}}
\toprule
$a$ & Weight $2a$ & $\dim S_{2a}$ & New cusp type & New $E_i$ found? \\
\midrule
8  & 16 & 1 & $\Delta\cdot E_4$              & No \\
9  & 18 & 1 & $\Delta\cdot E_6$              & No \\
10 & 20 & 1 & $\Delta\cdot E_8$              & No \\
11 & 22 & 1 & $\Delta\cdot E_{10}$           & No \\
12 & 24 & 2 & $\Delta\cdot E_{12},\;\Delta^2$ & No \\
\bottomrule
\end{tabular}
\end{table}


\section*{Appendix: Computational Data}
\addcontentsline{toc}{section}{Appendix: Computational Data}
\label{app:computations}

All computations were performed in exact rational arithmetic.
Prime evaluation matrices were built from the first $N$ primes and reduced via
rational row-reduction (RREF over $\Q$).
The source code is available at the project repository.

% ------------------------------------------------------------------
\subsection*{A.1\quad Closed Form of $M_6$}
\label{app:M6-formula}

Fitting the $55\times 22$ rational linear system
(21 polynomial-weighted divisor-sum columns plus one $\tau$-column)
yields the exact formula
\[
  M_6(n) = \sum_{j=0}^{5} P_j(n)\,\sigma_{2j+1}(n)
            - \frac{17}{150\,450\,048\,000}\,\tau(n),
\]
with explicit polynomials $P_j\in\Q[n]$:
\begin{align*}
P_0(n)&=\frac{550499}{4541644800}-\frac{153617}{371589120}\,n+\frac{2159}{6635520}\,n^2-\frac{67}{737280}\,n^3+\frac{11}{1105920}\,n^4-\frac{1}{2764800}\,n^5\\[4pt]
P_1(n)&=\frac{153617}{743178240}-\frac{2159}{8847360}\,n+\frac{67}{819200}\,n^2-\frac{11}{1105920}\,n^3+\frac{1}{2580480}\,n^4\\[4pt]
P_2(n)&=\frac{2159}{88473600}-\frac{67}{4915200}\,n+\frac{11}{5160960}\,n^2-\frac{1}{10321920}\,n^3\\[4pt]
P_3(n)&=\frac{67}{123863040}-\frac{11}{74317824}\,n+\frac{1}{111476736}\,n^2\\[4pt]
P_4(n)&=\frac{11}{3715891200}-\frac{1}{3096576000}\,n\\[4pt]
P_5(n)&=\frac{1}{268995133440}
\end{align*}
These 22 coefficients were verified against direct partition enumeration for
$n=1,\ldots,30$ with no mismatches (see repository script
\texttt{Paper/compute\_m6\_coefficients.jl}).
The key datum is the $\tau$-coefficient
\[
  c_\tau = -\frac{17}{150\,450\,048\,000},
  \qquad
  150\,450\,048\,000 = 2^{10}\times 3^5\times 5^3\times 7\times 691.
\]
The denominator prime $691 = \mathrm{num}(B_{12})$ is the Bernoulli prime
indexing the weight-12 Eisenstein series.

% ------------------------------------------------------------------
\subsection*{A.2\quad Cusp Structure of $M_7$ and Value of $\beta$}
\label{app:M7-formula}

Despite $\dim S_{14}(\SL_2(\Z))=0$, the closed form of $M_7(n)$ involves
$n\cdot\tau(n)$.
This follows from the quasimodular decomposition of $U_7(q)$: the depth-1 component
lies in $E_2\cdot M_{12}$, and the identity $D\Delta = E_2\Delta$ (proved from Ramanujan's differential equations
in the companion paper) shows that
$[q^n](E_2\Delta)=n\cdot\tau(n)$ exactly.
There is no independent $\tau(n)$ term: the $\tau(p)$ direction in $\mathcal{O}_d$
comes entirely from $M_6$ (via $c_\tau\ne 0$).

The exact closed form is
\[
  M_7(n) = \sum_{j=0}^{6} Q_j(n)\,\sigma_{2j+1}(n) + \beta\,n\cdot\tau(n),
\]
with $\beta\ne 0$.
The exact rational value of $\beta$ is the coefficient of $E_2\Delta$ in the
quasimodular decomposition of $U_7$, computed in exact arithmetic by the repository
script \texttt{Paper/extract\_e6\_e7.jl} (function \texttt{\_fit\_m7!}, column~30
of the augmented RREF system) and verified against direct partition enumeration for
$n=1,\dots,30$.

The structural significance is:
\begin{itemize}
  \item $\beta\ne 0$: $M_7$ contributes $p\cdot\tau(p)$ at primes, a direction absent
        from $M_6$ alone (which contributes only $\tau(p)$).
\end{itemize}
Together, the $\{M_6,M_7\}$ sub-system at primes carries both independent cusp
directions $\tau(p)$ and $p\cdot\tau(p)$, which is the source of the
cusp-cancellation mechanism underlying $E_6$.

% ------------------------------------------------------------------
\subsection*{A.3\quad Prime Evaluation Rank Data}
\label{app:rank-data}

\begin{table}[H]
\centering
\caption{Rank and nullity of prime evaluation matrices at $N=95$ primes,
         showing the pivot status of $M_6$ within the restricted basis.}
\label{tab:app-pivot}
\begin{tabular}{@{}lcccc@{}}
\toprule
Matrix & Rows & Cols & Rank & Nullity \\
\midrule
$a_{\max}=5,\;d=2$ & 95 & 15 & 11 & 4 \\
$a_{\max}=6,\;d=2$ & 95 & 18 & 14 & 4 \\
$a_{\max}=5,\;d=3$ & 95 & 20 & 12 & 8 \\
$a_{\max}=6,\;d=3$ & 95 & 24 & 16 & 8 \\
\bottomrule
\end{tabular}
\end{table}

In each case the nullity does not increase when $M_6$ is adjoined,
confirming that all $M_6$-columns are pivot columns.

% ------------------------------------------------------------------
\subsection*{A.4\quad Nullity Gains When $M_7$ Is Adjoined}
\label{app:nullity-m7}

\begin{table}[H]
\centering
\caption{Nullity comparison at $N=95$ primes.
         The gain first appears at $d=4$.}
\label{tab:app-nullity-m7}
\begin{tabular}{@{}ccccc@{}}
\toprule
$d$ & Nullity $\{M_1\text{--}M_5\}$ &
      Nullity $\{M_1\text{--}M_6\}$ &
      Nullity $\{M_1\text{--}M_7\}$ & Gain \\
\midrule
2 & 4  & 4  & 4  & 0 \\
3 & 8  & 8  & 8  & 0 \\
4 & 12 & 12 & 13 & $+1$ \\
5 & 16 & 16 & 18 & $+2$ \\
6 & 20 & 20 & 23 & $+3$ \\
\bottomrule
\end{tabular}
\end{table}

% ------------------------------------------------------------------
\subsection*{A.5\quad Explicit Formula for $E_5$}
\label{app:E5-formula}

After normalization to primitive integer coefficients,
the unique new prime-vanishing direction at degree~2 in the basis
$\{M_1,\dots,M_5\}$ is
\begin{align*}
E_5(n) ={}
  & (-450450+675675n-225225n^2)\,M_1(n) \\
  &+(960960n-120120n^2)\,M_2(n) \\
  &+(2534912n-166016n^2)\,M_3(n) \\
  &+(7999488n-322560n^2)\,M_4(n) \\
  &+258048000\,M_5(n).
\end{align*}
The constant $M_5$-coefficient is
$258048000 = 2^{11}\times 3^2\times 5^3\times 7\times 127$.

% ------------------------------------------------------------------
\subsection*{A.6\quad Explicit Formula for $E_6$}
\label{app:E6-formula}

The unique new prime-vanishing direction at degree~4 in the basis
$\{M_1,\dots,M_7\}$, normalized analogously, is
\begin{align*}
E_6(n) ={}
  &(105367732470 - 277943208789n + 289613157079n^2 \\
  &\quad{}-145583618619n^3+28545937859n^4)\,M_1(n)\\
  &+(-51324522800n^3+4292382480n^4)\,M_2(n)\\
  &+(-40313554176n^3+1776519808n^4)\,M_3(n)\\
  &+(-32888346624n^3+770273280n^4)\,M_4(n)\\
  &+(-7741440000n^3-154828800n^4)\,M_5(n)\\
  &+(-483548921856000+37196070912000n)\,M_6(n)\\
  &+892705701888000\,M_7(n).
\end{align*}
The leading $M_7$-coefficient is
$892705701888000 = 2^{10}\times 3^5\times 5^3\times 7\times 691\times 1303$.
The factor $691$ connects $E_6$ arithmetically to the weight-12 modular structure of
$M_6$.

% ------------------------------------------------------------------
\subsection*{A.7\quad Extended Search to $a_{\max}=12$}
\label{app:extended-search}

\begin{table}[H]
\centering
\caption{Extended search to $a_{\max}=12$.
         No independent $E_i$ was found beyond $E_6$.}
\label{tab:app-extended}
\begin{tabular}{@{}ccccc@{}}
\toprule
$a$ & Weight $2a$ & $\dim S_{2a}$ & New cusp type & New $E_i$ found? \\
\midrule
8  & 16 & 1 & $\Delta\cdot E_4$              & No \\
9  & 18 & 1 & $\Delta\cdot E_6$              & No \\
10 & 20 & 1 & $\Delta\cdot E_8$              & No \\
11 & 22 & 1 & $\Delta\cdot E_{10}$           & No \\
12 & 24 & 2 & $\Delta\cdot E_{12},\;\Delta^2$ & No \\
\bottomrule
\end{tabular}
\end{table}


\section*{Appendix: Computational Data}
\addcontentsline{toc}{section}{Appendix: Computational Data}
\label{app:computations}

All computations were performed in exact rational arithmetic.
Prime evaluation matrices were built from the first $N$ primes and reduced via
rational row-reduction (RREF over $\Q$).
The source code is available at the project repository.

% ------------------------------------------------------------------
\subsection*{A.1\quad Closed Form of $M_6$}
\label{app:M6-formula}

Fitting the $55\times 22$ rational linear system
(21 polynomial-weighted divisor-sum columns plus one $\tau$-column)
yields the exact formula
\[
  M_6(n) = \sum_{j=0}^{5} P_j(n)\,\sigma_{2j+1}(n)
            - \frac{17}{150\,450\,048\,000}\,\tau(n),
\]
with explicit polynomials $P_j\in\Q[n]$:
\begin{align*}
P_0(n)&=\frac{550499}{4541644800}-\frac{153617}{371589120}\,n+\frac{2159}{6635520}\,n^2-\frac{67}{737280}\,n^3+\frac{11}{1105920}\,n^4-\frac{1}{2764800}\,n^5\\[4pt]
P_1(n)&=\frac{153617}{743178240}-\frac{2159}{8847360}\,n+\frac{67}{819200}\,n^2-\frac{11}{1105920}\,n^3+\frac{1}{2580480}\,n^4\\[4pt]
P_2(n)&=\frac{2159}{88473600}-\frac{67}{4915200}\,n+\frac{11}{5160960}\,n^2-\frac{1}{10321920}\,n^3\\[4pt]
P_3(n)&=\frac{67}{123863040}-\frac{11}{74317824}\,n+\frac{1}{111476736}\,n^2\\[4pt]
P_4(n)&=\frac{11}{3715891200}-\frac{1}{3096576000}\,n\\[4pt]
P_5(n)&=\frac{1}{268995133440}
\end{align*}
These 22 coefficients were verified against direct partition enumeration for
$n=1,\ldots,30$ with no mismatches (see repository script
\texttt{Paper/compute\_m6\_coefficients.jl}).
The key datum is the $\tau$-coefficient
\[
  c_\tau = -\frac{17}{150\,450\,048\,000},
  \qquad
  150\,450\,048\,000 = 2^{10}\times 3^5\times 5^3\times 7\times 691.
\]
The denominator prime $691 = \mathrm{num}(B_{12})$ is the Bernoulli prime
indexing the weight-12 Eisenstein series.

% ------------------------------------------------------------------
\subsection*{A.2\quad Cusp Structure of $M_7$ and Value of $\beta$}
\label{app:M7-formula}

Despite $\dim S_{14}(\SL_2(\Z))=0$, the closed form of $M_7(n)$ involves
$n\cdot\tau(n)$.
This follows from the quasimodular decomposition of $U_7(q)$: the depth-1 component
lies in $E_2\cdot M_{12}$, and the identity $D\Delta = E_2\Delta$ (proved from Ramanujan's differential equations
in the companion paper) shows that
$[q^n](E_2\Delta)=n\cdot\tau(n)$ exactly.
There is no independent $\tau(n)$ term: the $\tau(p)$ direction in $\mathcal{O}_d$
comes entirely from $M_6$ (via $c_\tau\ne 0$).

The exact closed form is
\[
  M_7(n) = \sum_{j=0}^{6} Q_j(n)\,\sigma_{2j+1}(n) + \beta\,n\cdot\tau(n),
\]
with $\beta\ne 0$.
The exact rational value of $\beta$ is the coefficient of $E_2\Delta$ in the
quasimodular decomposition of $U_7$, computed in exact arithmetic by the repository
script \texttt{Paper/extract\_e6\_e7.jl} (function \texttt{\_fit\_m7!}, column~30
of the augmented RREF system) and verified against direct partition enumeration for
$n=1,\dots,30$.

The structural significance is:
\begin{itemize}
  \item $\beta\ne 0$: $M_7$ contributes $p\cdot\tau(p)$ at primes, a direction absent
        from $M_6$ alone (which contributes only $\tau(p)$).
\end{itemize}
Together, the $\{M_6,M_7\}$ sub-system at primes carries both independent cusp
directions $\tau(p)$ and $p\cdot\tau(p)$, which is the source of the
cusp-cancellation mechanism underlying $E_6$.

% ------------------------------------------------------------------
\subsection*{A.3\quad Prime Evaluation Rank Data}
\label{app:rank-data}

\begin{table}[H]
\centering
\caption{Rank and nullity of prime evaluation matrices at $N=95$ primes,
         showing the pivot status of $M_6$ within the restricted basis.}
\label{tab:app-pivot}
\begin{tabular}{@{}lcccc@{}}
\toprule
Matrix & Rows & Cols & Rank & Nullity \\
\midrule
$a_{\max}=5,\;d=2$ & 95 & 15 & 11 & 4 \\
$a_{\max}=6,\;d=2$ & 95 & 18 & 14 & 4 \\
$a_{\max}=5,\;d=3$ & 95 & 20 & 12 & 8 \\
$a_{\max}=6,\;d=3$ & 95 & 24 & 16 & 8 \\
\bottomrule
\end{tabular}
\end{table}

In each case the nullity does not increase when $M_6$ is adjoined,
confirming that all $M_6$-columns are pivot columns.

% ------------------------------------------------------------------
\subsection*{A.4\quad Nullity Gains When $M_7$ Is Adjoined}
\label{app:nullity-m7}

\begin{table}[H]
\centering
\caption{Nullity comparison at $N=95$ primes.
         The gain first appears at $d=4$.}
\label{tab:app-nullity-m7}
\begin{tabular}{@{}ccccc@{}}
\toprule
$d$ & Nullity $\{M_1\text{--}M_5\}$ &
      Nullity $\{M_1\text{--}M_6\}$ &
      Nullity $\{M_1\text{--}M_7\}$ & Gain \\
\midrule
2 & 4  & 4  & 4  & 0 \\
3 & 8  & 8  & 8  & 0 \\
4 & 12 & 12 & 13 & $+1$ \\
5 & 16 & 16 & 18 & $+2$ \\
6 & 20 & 20 & 23 & $+3$ \\
\bottomrule
\end{tabular}
\end{table}

% ------------------------------------------------------------------
\subsection*{A.5\quad Explicit Formula for $E_5$}
\label{app:E5-formula}

After normalization to primitive integer coefficients,
the unique new prime-vanishing direction at degree~2 in the basis
$\{M_1,\dots,M_5\}$ is
\begin{align*}
E_5(n) ={}
  & (-450450+675675n-225225n^2)\,M_1(n) \\
  &+(960960n-120120n^2)\,M_2(n) \\
  &+(2534912n-166016n^2)\,M_3(n) \\
  &+(7999488n-322560n^2)\,M_4(n) \\
  &+258048000\,M_5(n).
\end{align*}
The constant $M_5$-coefficient is
$258048000 = 2^{11}\times 3^2\times 5^3\times 7\times 127$.

% ------------------------------------------------------------------
\subsection*{A.6\quad Explicit Formula for $E_6$}
\label{app:E6-formula}

The unique new prime-vanishing direction at degree~4 in the basis
$\{M_1,\dots,M_7\}$, normalized analogously, is
\begin{align*}
E_6(n) ={}
  &(105367732470 - 277943208789n + 289613157079n^2 \\
  &\quad{}-145583618619n^3+28545937859n^4)\,M_1(n)\\
  &+(-51324522800n^3+4292382480n^4)\,M_2(n)\\
  &+(-40313554176n^3+1776519808n^4)\,M_3(n)\\
  &+(-32888346624n^3+770273280n^4)\,M_4(n)\\
  &+(-7741440000n^3-154828800n^4)\,M_5(n)\\
  &+(-483548921856000+37196070912000n)\,M_6(n)\\
  &+892705701888000\,M_7(n).
\end{align*}
The leading $M_7$-coefficient is
$892705701888000 = 2^{10}\times 3^5\times 5^3\times 7\times 691\times 1303$.
The factor $691$ connects $E_6$ arithmetically to the weight-12 modular structure of
$M_6$.

% ------------------------------------------------------------------
\subsection*{A.7\quad Extended Search to $a_{\max}=12$}
\label{app:extended-search}

\begin{table}[H]
\centering
\caption{Extended search to $a_{\max}=12$.
         No independent $E_i$ was found beyond $E_6$.}
\label{tab:app-extended}
\begin{tabular}{@{}ccccc@{}}
\toprule
$a$ & Weight $2a$ & $\dim S_{2a}$ & New cusp type & New $E_i$ found? \\
\midrule
8  & 16 & 1 & $\Delta\cdot E_4$              & No \\
9  & 18 & 1 & $\Delta\cdot E_6$              & No \\
10 & 20 & 1 & $\Delta\cdot E_8$              & No \\
11 & 22 & 1 & $\Delta\cdot E_{10}$           & No \\
12 & 24 & 2 & $\Delta\cdot E_{12},\;\Delta^2$ & No \\
\bottomrule
\end{tabular}
\end{table}


\end{document}
